\documentclass[11pt,a4paper]{article}
\usepackage{amsmath,amssymb}
\newcommand{\Tr}{\operatorname{Tr}}
\begin{document}
\title{The global instantaneous observational-entropy maximum for binary POVMs}
\author{}
\date{}
\maketitle
\paragraph{Setting and claim.} Let the Hilbert space have finite dimension $d\geq2$, and let $\rho=|\psi\rangle\langle\psi|$ evolve differentiably and unitarily at the instant considered, with generator $H/\hbar$. All measurements are frozen when taking the time derivative. Use natural logarithms and full-space trace volumes. Set $s=\Delta H/\hbar$. Among all binary POVMs $\{E,I-E\}$, including nonprojective effects, the maximum instantaneous observational-entropy rate is
\[R_{\max,2}=s C_{d,2},\qquad C_{d,2}=2\sqrt{p_*(1-p_*)}\log\frac{(d-1)p_*}{1-p_*}.\]
Here $p_*$ is the unique solution in $(1/2,1)$ of
\[(2p_*-1)\log\frac{(d-1)p_*}{1-p_*}=2.\]
If $s=0$, all instantaneous rates vanish. For $s>0$, define $|0\rangle=|\psi\rangle$ and $|1\rangle=i(H-\langle H\rangle)|\psi\rangle/\Delta H$. Then an attaining measurement is
\[E_*=|u\rangle\langle u|,\qquad |u\rangle=\sqrt{p_*}|0\rangle+\sqrt{1-p_*}|1\rangle.\]
The other effect is $I-E_*$. The proof below establishes global maximality by domination of every binary effect, and separately supplies an explicit KKT certificate. KKT stationarity alone would not establish global maximality.
\paragraph{1. Normalize the tangent and reduce the variables.} Divide the rate by $s$, so that the normalized derivative is
\[\dot\rho_0=-(|0\rangle\langle1|+|1\rangle\langle0|).\]
For a general effect $0\preceq E\preceq I$, write $a=\langle0|E|0\rangle$, $c=\langle1|E|1\rangle$, $b=\operatorname{Re}\langle0|E|1\rangle$ and $V=\Tr E$. With $0<a<1$ and $0<V<d$, the dimensionless rate is
\[r=R/s=2bL(a,V),\qquad L(a,V)=\log\frac{a(d-V)}{(1-a)V}.\]
Relabelling the two outcomes sends $(a,V,b)$ to $(1-a,d-V,-b)$ and leaves $r$ invariant. Consequently we may assume $L\geq0$. Nonpositive rates are already dominated by the positive candidate, so it remains to treat $L>0$ and $b>0$.

The two principal positivity constraints imply
\[b^2\leq|\langle0|E|1\rangle|^2\leq\min\{ac,(1-a)(1-c)\}.\]
Write $V=a+c+w$, where $w\geq0$ is the trace on the orthogonal complement. The principal two-dimensional block is itself a positive contraction, so replacing the complement and its couplings by zero is feasible. Since $L(a,V)$ strictly decreases with $V$, this replacement cannot lower a positive rate. Replacing the off-diagonal element by the positive real saturation of its bound is also feasible for the two-dimensional block. Thus every positive binary rate is bounded by
\[2\sqrt{\min\{ac,(1-a)(1-c)\}}\,L(a,a+c).\]
This argument also covers $d=2$, where there is no complement.
\paragraph{2. Every effect is dominated by a rank-one projector.} There are two cases.

If $a+c\geq1$, the smaller determinant bound is $(1-a)(1-c)$. Decreasing $c$ to $1-a$ increases this bound and decreases the trace, hence increases the logarithmic factor. Both factors are nonnegative throughout the comparison. The resulting upper bound is achieved by the projector with diagonal entries $a,1-a$. In particular, positivity of the logarithm implies $ad>a+c\geq1$, so $a>1/d$.

If $0<a+c=t\leq1$, the smaller bound is $ac$. Saturating it gives a rank-one effect $tP$, with $a=tp$ and $0<p<1$. Its rate is
\[r(t,p)=2t\sqrt{p(1-p)}\log\frac{p(d-t)}{1-tp}.\]
The logarithm is positive exactly when $dp>1$. At fixed such $p$, its derivative satisfies
\[\frac{\partial}{\partial t}\left[t\log\frac{p(d-t)}{1-tp}\right]=\log\frac{p(d-t)}{1-tp}+\frac{t(dp-1)}{(d-t)(1-tp)}>0.\]
Therefore increasing $t$ to one improves the rate and yields a rank-one projector. This proves a global comparison with a projective measurement, not merely a stationarity condition.
\paragraph{3. Solve the remaining scalar problem.} The preceding steps reduce the full binary problem to maximizing
\[f_d(p)=2\sqrt{p(1-p)}\,L_d(p),\qquad L_d(p)=\log\frac{(d-1)p}{1-p},\quad 1/d\leq p\leq1.\]
The value at both endpoints is zero, using the continuous limit at $p=1$. In the interior,
\[f_d'(p)=\frac{2-(2p-1)L_d(p)}{\sqrt{p(1-p)}}.\]
The derivative is positive on $[1/d,1/2]$ wherever this interval has interior. On $(1/2,1)$, the function $(2p-1)L_d(p)$ increases strictly from zero to infinity, because
\[\frac{d}{dp}\big[(2p-1)L_d(p)\big]=2L_d(p)+\frac{2p-1}{p(1-p)}>0.\]
There is consequently exactly one stationary point $p_*$ on this interval, and it is the global maximum. This proves the claim.
\paragraph{4. Explicit KKT certificate.} Continue with unit normalized speed. For outcome probabilities $p_i=\Tr(\rho E_i)$, derivatives $\dot p_i=\Tr(\dot\rho_0 E_i)$ and volumes $V_i=\Tr E_i$, the effect gradients are
\[G_i=-\log(p_i/V_i)\dot\rho_0-\frac{\dot p_i}{p_i}\rho+\frac{\dot p_i}{V_i}I.\]
At the projector candidate, abbreviate $p=p_*$, $q=1-p$, $z=\sqrt{pq}$, $L=L_d(p)$ and $\beta=2zd/(d-1)$. For $P=|u\rangle\langle u|$ one obtains
\[D:=G_1-G_2=-L\dot\rho_0+\frac{2}{z}\rho-\beta I.\]
Let $|w\rangle=-\sqrt q|0\rangle+\sqrt p|1\rangle$. The off-diagonal matrix element is
\[\langle w|D|u\rangle=(2p-1)L-2=0.\]
Thus $D$ is diagonal in the decomposition into $u,w$ and the remaining complement, with eigenvalues
\[\lambda_+=\frac{2p}{z(2p-1)}-\beta>0,\qquad \lambda_-=-\frac{2q}{z(2p-1)}-\beta<0.\]
On the remaining complement its eigenvalue is $-\beta<0$. To check $\lambda_+>0$, use $1/(2p-1)>1$ and $qd/(d-1)<1$ for $p>1/2$ and $d\geq2$. Define the common Hermitian multiplier
\[\Lambda=G_2+\lambda_+P.\]
Then the binary POVM satisfies the semidefinite KKT conditions
\[\Lambda-G_i\succeq0,\qquad(\Lambda-G_i)E_i=0,\qquad E_1+E_2=I.\]
The KKT certificate verifies first-order optimality; the global guarantee comes from the domination argument in steps 1--3.
\paragraph{Boundary cases and scope.} If a fixed effect has zero probability at the instant considered, its probability along a smooth finite-dimensional unitary orbit is $O(\tau^2)$; its entropy contribution is $O(\tau^2|\log|\tau||)$ and has zero first derivative. Thus the excluded $a=0,1$ boundaries do not improve the result. A zero-volume effect is the zero operator and can be omitted. The result concerns a pure state at one instant, not mixed states, a common POVM for an ensemble, or a sustained slope over an interval. Time-dependent Hamiltonians are allowed at a differentiable instant; an ideal delta kick itself is excluded.

Proportional splitting of $I-E_*$ produces any prescribed number of nonzero outcomes without changing the rate, so this theorem supplies an attainable lower bound for every outcome count at least two. It does not show that arbitrary higher-outcome POVMs satisfy the same upper bound. Such a bound, or a valid merging argument, is still required before attempting induction.
\paragraph{Numerical cross-checks.} The script \texttt{calculations/code/check\_binary\_obsen\_global\_bound.py} checks the scalar equation, KKT slacks and complementarity, and random feasible binary effects at $d=2,3,21,256,4096$. These checks supplement rather than replace the analytic proof.
\end{document}
